\chapter{Εισαγωγή}

Η εκτίμηση ηλεκτρομηχανικών ταλαντώσεων από \textit{ambient} δεδομένα αποτελεί χρήσιμο εργαλείο για τη μελέτη της δυναμικής συμπεριφοράς Συστημάτων Ηλεκτρικής Ενέργειας χωρίς την ανάγκη έντονων διαταραχών. Σε αντίθεση με την \textit{ringdown} ανάλυση, όπου η πληροφορία εξάγεται από ορατές αποκρίσεις μετά από συμβάν, στο \textit{ambient} πλαίσιο η modal πληροφορία πρέπει να εξαχθεί από ασθενέστερες ταλαντωτικές συνιστώσες.

Η μελέτη επικεντρώνεται στην παραγωγή modal εκτιμήσεων από προεπεξεργασμένες χρονοσειρές τάσης και ρεύματος ($V$, $I$) για το \textit{IEEE 39-Bus System}~\cite{digsilent39}. Παράλληλα, αξιοποιείται ξεχωριστή modal εξαγωγή από το \textit{PowerFactory}, ώστε να προκύπτουν modes αναφοράς ειδικά για το ίδιο \textit{ambient} dataset. Με τον τρόπο αυτό, η αξιολόγηση συνδέεται με το modal περιεχόμενο του ίδιου του σεναρίου προσομοίωσης.
